\subsection{Evaluation setting}
\label{sec:setup}

\textbf{Research questions.}
\textbf{RQ1} How accurate is \nscg compared with PyCG-2021 and PyCG-2026?
\textbf{RQ2} What does each component of \nscg contribute, and how much does the
model matter?
\textbf{RQ3} Which edges does each approach get wrong, and why?
\textbf{RQ4} Do the LLM's answers depend on memorised names, and are they stable
across runs?
\textbf{RQ5} Do hybrids of PyCG and \nscg combine their strengths?

\textbf{Subjects.} We use the two benchmarks of the PyCG evaluation repository~\cite{pycgeval2026}
(commit \texttt{8a2e52c}), because they are the only Python call-graph
benchmarks with executable baselines and carefully defined ground truths.
(i) The \emph{micro-benchmark}: the 622 cases up to Python~3.12 (the cohort the
evaluation uses to compare tools), small programs of 34 feature categories
(classes, decorators, dynamic features, imports, standard library, \dots), each
with a ground truth validated against a dynamic trace. (ii) Four packages of the
\emph{macro-benchmark} (Table~\ref{tab:subjects}), whose ground truths are
call-site ledgers written by hand, checked against the traces of the packages'
test suites, and audited four times. Given a limited budget and the course's
guidance of three to four real projects, we selected them to cover both regimes
of the baseline: two where PyCG-2026 is strong (autojump, a small script-like
utility; jmespath, an interpreter with visitor dispatch through
\texttt{getattr}) and two where it is weak (fabric, which subclasses and is
called back by the Invoke framework; funcy, a higher-order functional library),
two from the ICSE~2021 set and two from later batches.

\begin{table}[t]
  \caption{The four packages (sources at the pinned commits of the benchmark).
  Edges in/out: ground-truth edges inside / outside the package. Sites: call
  sites found by the \nscg front-end.}
  \label{tab:subjects}
  \small
  \setlength{\tabcolsep}{3.2pt}
  \begin{tabular}{l l r r r r r r}
\toprule
Package & batch & LOC & mod. & nodes & edges in/out & sites & PyCG F1\\
\midrule
autojump & ICSE'21 & 488 & 3 & 46 & 35/100 & 139 & 0.95 \\
fabric & ICSE'21 & 2,303 & 12 & 93 & 99/252 & 377 & 0.71 \\
jmespath & 1 & 1,675 & 8 & 190 & 310/120 & 470 & 0.99 \\
funcy & 5 & 2,392 & 16 & 305 & 511/475 & 987 & 0.65 \\
\bottomrule
\end{tabular}

\end{table}

\textbf{Criteria.} Edge precision, recall and $F_1$, computed by the official
scorer (\texttt{evaluation/common/score.py}, raw names) separately for callees
inside the package (the headline, as in the PyCG evaluation) and outside it; for
the micro-benchmark also the number of cases whose graph equals the ground truth
exactly. A run that fails counts every edge of its program as missed.

\textbf{Setup.} All runs use one machine (Ubuntu 22.04, 128 cores, 755\,GB
RAM) and Python~3.14.3. The baselines run through the official runner with a
30-minute and 16\,GiB limit per program (PyCG-2021 timed out on funcy, as in
the PyCG evaluation). \nscg uses Claude Haiku~5.5 (and Sonnet~5.5 for RQ2)
through \texttt{claude -p} 2.1.321, 16 requests in parallel, default sampling
(the CLI exposes no temperature), one sample per request; RQ4 repeats the
macro runs with independent samples. The same subjects, entry points and
package roots are given to every tool, and no tool sees the ground truth. The
skill text of the prompts was written from the benchmark's documented
conventions (the same information the PyCG authors had), and the prompts were
developed on a handful of micro cases only; the four packages were not used to
tune the prompts before the V2 run. The total LLM cost of every request of the
study is US\$21.15 (Table~\ref{tab:cost}), of which the
Sonnet runs are three quarters. Money was not the binding constraint; the time
frame of the project (one week of wall-clock time, including the 30-minute
limits of the baselines) and the course's guidance limited the macro
evaluation to four packages, and we chose the cheapest model as the default so
that every variant could be run on all subjects.
